本稿は, 水文頻度解析で用いられる標準最小二乗規準 (SLSC) について, 標本サイズ依存性と標準化の影響を再検討した.
6分布を対象としたモンテカルロ計算により, SLSC は標本サイズの増加とともに低下する傾向を示した.
また, \(X\) 空間と \(S\) 空間を区別し, 非線形標準化では SLSC が標準化関数に依存する距離評価となることを示した.
特に LN3 の \(S\) 空間評価は小さい値を与えやすく, 分布間比較に影響しうる.
\rev{さらに, SLSC+ジャックナイフ法と AIC の採択結果も比較した.}
SLSC の解釈には, 固定閾値だけでなく, 標準変量の定義と距離尺度の違いを明示する必要がある.
